Dengue surveillance in Bangladesh operates at the intersection of a recurrent public-health problem and a reporting system that must support rapid operational decisions. The country's dengue evidence base includes hospital and clinical reporting, outbreak communications, spatial analyses, and predictive work, while the national surveillance architecture remains dependent on the quality, timeliness, and interpretability of the observations that reach public reporting interfaces \citep{ref1,ref6}. This distinction matters because a value that is publicly visible is not automatically equivalent to a complete measure of infection, a population denominator, or a stable representation of the underlying surveillance system.

Recent literature demonstrates substantial analytic activity using Bangladesh dengue observations. National analyses have examined temporal and geographic patterns across 2019--2023 \citep{ref2}, and district-level hotspot and cluster methods have been applied to the 2023 outbreak \citep{ref3}. A 2026 retrospective ecological study extends this broad precedent across all 64 districts using reported hospital admissions \citep{ref4}. These studies establish that national, district, and spatial dengue analyses are already represented in the literature. They also make the provenance and interpretation of shared surveillance observations consequential: the scientific question is not only what a dashboard value appears to show, but also which reporting universe, time window, geography, and revision state it represents.

System-level assessments provide an important reason to examine these properties explicitly. A recent narrative review describes Bangladesh's dengue surveillance as involving passive and active modalities, with passive and hospital-based reporting offering wide coverage but remaining vulnerable to underreporting and delayed notification \citep{ref1}. The same review identifies variability in data quality, timeliness, interoperability, and terminology, and notes limited standardization in the explicit reporting of case or outbreak definitions \citep{ref1}. The World Health Organization's Bangladesh country assessment likewise treats communicable-disease surveillance as an architecture requiring mapping, performance evaluation, and opportunities for integration and improvement, with dengue included among the priority diseases considered \citep{ref6}. An official 2023 situation report further describes hospital-based surveillance, regular information from hospitals and Upazila and district facilities, and daily dissemination through the Health Emergency Operations Center \citep{ref7}. Together, these sources support a layered interpretation: public dashboard observations are outputs of a surveillance and reporting environment, not transparent substitutes for that environment.

The public DGHS Dengue Dynamic Dashboard is therefore an important research-data interface. Its accessibility can support reproducible observation of historical dashboard states, but accessibility alone does not establish that a response can be reconstructed, that displayed fields are semantically interchangeable, or that geographic labels identify patient residence rather than a facility or administrative reporting universe. A publicly available non-peer-reviewed reproducible analysis repository already describes a national analysis based on DGHS HEOC dashboard-derived aggregate surveillance data \citep{ref5}. That repository is relevant grey-literature prior art and further prohibits framing reproducibility of DGHS-derived dengue data as a first-in-field contribution. Its presence also reinforces the need to state precisely what a new study audits: the source interface, its exposed observations, and the evidentiary limits surrounding their interpretation.

ResearchDengu addresses this narrower methodological problem. The study evaluates whether selected historical states of the public DGHS dashboard can be reconstructed and preserved as research artifacts, and whether their source-labelled fields can be interpreted without silently converting them into population-level measures. The primary object is the dashboard as a public surveillance-data interface rather than Bangladesh dengue burden, transmission, climate association, spatial risk, or prediction. Restricted descriptive dengue observations are retained to show the structure and use of the interface, but they are not used to estimate incidence, prevalence, risk, causal effects, or the completeness of the national surveillance system.

The study has four linked aims. First, it documents the provenance of archived dashboard responses through date-specific acquisition records and cryptographic hashes. Second, it evaluates whether the observed HTML, chart, control, and table structures support reproducible reconstruction of selected dashboard states. Third, it preserves raw labels and tests semantic questions that affect comparability, including reporting windows, category definitions, and geographic meaning. Fourth, it records internal reconciliation limitations and unresolved revision or stability questions without correcting, imputing, or biologically interpreting them. The resulting contribution is consequently a provenance-preserving observation-level reproducibility and data-quality audit.

This framing is deliberately bounded by the existing literature. ResearchDengu does not claim to be the first nationwide dengue analysis, the first reproducible Bangladesh dengue dataset, or the first surveillance-quality study in Bangladesh \citep{ref1,ref2,ref3,ref4,ref5,ref6,ref7}. Instead, it complements national and district epidemiological studies by making the data-interface layer inspectable: what was requested, what was returned, what was displayed, how it was archived, and which claims remain unsupported. The Methods and Results that follow therefore distinguish dashboard-level reproducibility from underlying surveillance validity and describe all retained values using the source terminology and limitations recorded in the project's evidence contracts.
